% Part: first-order-logic
% Chapter: completeness
% Section: lindenbaums-lemma

\documentclass[../../../include/open-logic-section]{subfiles}

\begin{document}

\iftag{FOL}
      {\olfileid{fol}{com}{lin}}
      {\olfileid{pl}{com}{lin}}

\olsection{لمّة ليندنبوم}

\begin{explain}
المقصود بهذه اللمّة أن كل مجموعة متسقة من ال!!{sentence}s لها امتداد
من الجمل يجمع إلى الاتساق كونه !!{complete}. وطريق البرهان أن نزيد
!!{sentence} واحدة في كل مرحلة، على شرط ألا يزول الاتساق. نرتب
الزيادات بحيث كل $!A$ يلحق إما $!A$ وإما $\lnot !A$ عند مرحلة ما.
فإذا أخذنا اتحاد المراحل كلها، وجدنا فيه $!A$ أو نفيها~$\lnot !A$،
فكان كاملًا. ويبقى متسقًا لأن كل مرحلة منه مصونة من التناقض.
\end{explain}

\begin{lem}[لمّة ليندنبوم]
\ollabel{lem:lindenbaum} يمكن توسيع كل مجموعة متسقة~$\Gamma$ في
لغة~$\Lang{L}$ إلى مجموعة !!a{complete} ومتسقة~$\Gamma^*$.
\end{lem}

\begin{proof}
ليكن الفرض اتساق $\Gamma$. نرتب $!A_0$ و$!A_1$ و\dots{} في قائمة
تستوعب ال!!{sentence}s في~$\Lang L$، ثم نضع $\Gamma_0 = \Gamma$، و
\[
\Gamma_{n+1} =
\begin{cases}
\Gamma_n \cup \{ !A_n \} & \textrm{متى كانت $\Gamma_n \cup \{!A_n\}$
  متسقة؛} \\
\Gamma_n \cup \{ \lnot !A_n \} & \textrm{وإلا.}
\end{cases}
\]
ولتكن $\Gamma^* = \bigcup_{n \geq 0} \Gamma_n$.

نثبت استقراءً اتساق كل $\Gamma_n$. فأما $\Gamma_0$ فمتسقة بحكم
تعريفها. وإذا اخترنا $\Gamma_{n+1} = \Gamma_n \cup \{!A_n\}$، فإنما
اخترناها لاتساقها. وإلا كان $\Gamma_{n+1} = \Gamma_n \cup \{\lnot !A_n\}$،
فلنثبت اتساق $\Gamma_n \cup \{\lnot !A_n\}$. لو لم تكن متسقة، لكانت
\emph{كل من} $\Gamma_n \cup \{!A_n\}$ و$\Gamma_n \cup \{\lnot !A_n\}$
غير متسقة؛ فيلزم عدم اتساق $\Gamma_n$ بمقتضى
\iftag{FOL}{%
  \tagrefs{prfAX/{fol:axd:prv:prop:provability-exhaustive},
    prfSC/{fol:seq:prv:prop:provability-exhaustive},
    prfND/{fol:ntd:prv:prop:provability-exhaustive},
    prfTab/{fol:tab:prv:prop:provability-exhaustive}}}{%
  \tagrefs{prfAX/{pl:axd:prv:prop:provability-exhaustive},
    prfSC/{pl:seq:prv:prop:provability-exhaustive},
    prfND/{pl:ntd:prv:prop:provability-exhaustive},
    prfTab/{pl:tab:prv:prop:provability-exhaustive}}},
خلافًا لفرض الاستقراء.

ونثبت أيضًا أنه لكل~$n$ ولكل $i < n$، يكون $\Gamma_i \subseteq \Gamma_n$،
باستقراء على~$n$. ففي حالة $n=0$ ليس ثمة $i < 0$، فتصدق الدعوى
خلوًا. وافترض الآن صحتها عند~$n$؛ نريد أن $i < n+1$ يقتضي
$\Gamma_i \subseteq \Gamma_{n+1}$. وبحسب البناء إما
$\Gamma_{n+1} = \Gamma_n \cup \{!A_n\}$ وإما $= \Gamma_n \cup \{\lnot
!A_n\}$؛ فيكون $\Gamma_n \subseteq \Gamma_{n+1}$. ومتى كان $i < n+1$،
كان $\Gamma_i \subseteq \Gamma_n$: بفرض الاستقراء إن كان $i < n$،
وبالانعكاسية $\Gamma_n \subseteq \Gamma_n$ إن كان $i = n$. ثم تفيد
تعدية~$\subseteq$ أن $\Gamma_i \subseteq \Gamma_{n+1}$.

بهذا يتبين اتساق $\Gamma^*$. خذ مجموعة منتهية
$\Gamma' \subseteq \Gamma^*$. فإن كان $\Gamma'=\emptyset$، فالاتساق
ظاهر. وإلا فكل $!B \in \Gamma'$ واقع في~$\Gamma_i$ عند فهرس~$i$.
اختر $n$ أكبر هذه الفهارس. وحيث إن
$\Gamma_i \subseteq \Gamma_n$ متى كان $i \le n$، فكل
$!B \in \Gamma'$ يحقق $!B \in \Gamma_n$، فيكون
$\Gamma' \subseteq \Gamma_n$ مع اتساق~$\Gamma_n$. فثبت اتساق كل
مجموعة منتهية $\Gamma' \subseteq \Gamma^*$. وتفيد \iftag{FOL}{%
  \tagrefs{prfAX/{fol:axd:ptn:prop:proves-compact},
    prfSC/{fol:seq:ptn:prop:proves-compact},
    prfND/{fol:ntd:ptn:prop:proves-compact},
    prfTab/{fol:tab:ptn:prop:proves-compact}}}{%
  \tagrefs{prfAX/{pl:axd:ptn:prop:proves-compact},
    prfSC/{pl:seq:ptn:prop:proves-compact},
    prfND/{pl:ntd:ptn:prop:proves-compact},
    prfTab/{pl:tab:ptn:prop:proves-compact}}}، تكون $\Gamma^*$ متسقة.

وأما الاكتمال، فكل !!{sentence} من $\Frm[L]$ لها موضع في قائمة بناء
~$\Gamma^*$. فإن كان $!A_n \notin \Gamma^*$، لم يمنع إلحاقها إلا
عدم اتساق $\Gamma_n \cup \{!A_n\}$؛ وحينئذ قد ألحقنا
$\lnot !A_n \in \Gamma^*$. فالمجموعة $\Gamma^*$ !!{complete}.
\end{proof}

\end{document}
